\par\addvspace{0.55\baselineskip}
\noindent\textbf{3.2. Equivalence relation defined by a group operating freely}\label{IV.sec.3.2}
\par\addvspace{0.35\baselineskip}

\begin{definition}{3.2.1}
\label{IV.3.2.1}
Let \(X\) be an object of~\(\mathcal{C}\) and \(H\) a \(\mathcal{C}\)-group operating
on~\(X\) (that is, equipped with a morphism of
\(\widehat{\mathcal{C}}\)-groups
\[
H\longrightarrow\Aut(X)\,).
\]
One says that \(H\) operates \emph{freely} on~\(X\) if the following equivalent
conditions are satisfied:
\begin{enumeratei}
\item For every \(S\in\Ob(\mathcal{C})\), the group \(H(S)\) operates freely
on~\(X(S)\);
\item The morphism of functors
\[
H\times X\longrightarrow X\times X
\]
defined set-theoretically by \((h,x)\mapsto(hx,x)\) is a monomorphism.
\end{enumeratei}
(In the preceding definition, one has supposed that the group \(H\) operated
``on the left'' on~\(X\). One evidently has an analogous notion in the case
where the group operates ``on the right'', that is, when one has been given
a morphism of groups from the group \(H^{\circ}\) opposite to~\(H\) into
\(\Aut(X)\).)
\end{definition}

If \(H\) operates freely on~\(X\), the image of \(H\times X\) by the morphism
of~(ii) is an equivalence relation in~\(X\), called the \emph{equivalence
relation defined by the action of~\(H\) on~\(X\)}. When the quotient of~\(X\)
by this equivalence relation exists, one denotes it by \(H\backslash X\)
(\(X/H\) when \(H\) operates on the right). It represents the following
covariant functor: if \(Z\) is an object of~\(\mathcal{C}\), one has
\[
\Hom(H\backslash X,Z)
=\{\text{morphisms from \(X\) into \(Z\) invariant under \(H\)}\}
\]
where the morphism \(f\colon X\to Z\) is said to be invariant under~\(H\) if
for every \(S\in\Ob(\mathcal{C})\), the corresponding morphism
\(X(S)\to Z(S)\) is invariant under the group \(H(S)\).

\begin{lemma}{3.2.2}
\label{IV.3.2.2}
Under the conditions of~3.2.1, let \(Y\) be a subobject of~\(X\). The
following conditions are equivalent:
\begin{enumeratei}
\item \(Y\) is stable under the equivalence relation defined by~\(H\)
(3.1.5);
\item For every \(S\in\Ob(\mathcal{C})\), the subset \(Y(S)\) of~\(X(S)\) is
stable under \(H(S)\);
\item There exists a morphism~\(f\), necessarily unique, making the following
diagram commutative
\[
\xymatrix@C=4.2em@R=2.4em{
H\times Y \ar[r]^{f} \ar[d]
  & Y \ar[d] \\
H\times X \ar[r]
  & X\,.
}
\]
\end{enumeratei}
Under these conditions, \(f\) defines a morphism of
\(\widehat{\mathcal{C}}\)-groups
\[
H\longrightarrow\Aut(Y)
\]
and the equivalence relation defined in~\(Y\) by this operation of~\(H\) is
none other than the equivalence relation induced in~\(Y\) by the equivalence
relation defined in~\(X\) by the action of~\(H\).
\end{lemma}

\begin{proof}
The proof is immediate. One evidently has an analogous statement for a
``right operation''. The operation of~\(H\) on~\(Y\) defined above will be
called the operation \emph{induced} in~\(Y\) by the given operation
of~\(H\) on~\(X\).
\end{proof}

Let us now consider the following situation: \(H\) and \(G\) are two
\(\mathcal{C}\)-groups and one has been given a morphism of groups
\[
u\colon H\longrightarrow G.
\]
Then \(H\) operates on~\(G\) by translations (one sets, set-theoretically,
\(hg=u(h)g\)) and operates there \emph{freely} if and only if \(u\) is a
\emph{monomorphism}. The quotient of~\(G\) by this operation of~\(H\) is
denoted, when it exists, by \(H\backslash G\). One likewise defines a right
operation of~\(H\) on~\(G\) and a quotient \(G/H\). These quotients are
functorial with respect to the groups in question; precisely, one has the
following lemma, stated for the right quotients:

\begin{lemma}{3.2.3}
\label{IV.3.2.3}
Let \(u\colon H\to G\) and \(u'\colon H'\to G'\) be two monomorphisms of
\(\mathcal{C}\)-groups. Suppose given a morphism of \(\mathcal{C}\)-groups
\[
f\colon G\longrightarrow G'.
\]
% typo: condition -> conditions
The following conditions are equivalent:
\begin{enumeratei}
\item \(f\) is compatible with the equivalence relations defined in \(G\)
and~\(G'\) by \(H\) and~\(H'\).
% typo: équivalences -> équivalence
\item For every \(S\in\Ob(\mathcal{C})\), one has
\(fu(H(S))\subset u'(H'(S))\).
\item There exists a morphism \(g\colon H\to H'\), necessarily unique and
multiplicative, such that the following diagram is commutative
\[
\xymatrix@C=4.2em@R=2.4em{
H \ar[r]^{g} \ar[d]_{u}
  & H' \ar[d]^{u'} \\
G \ar[r]_{f}
  & G'\,.
}
\]
\end{enumeratei}
Under these conditions, if the quotients \(G/H\) and \(G'/H'\) exist, there
exists a unique morphism~\(\overline{f}\) making the following diagram commutative
\[
\xymatrix@C=4.2em@R=2.4em{
G \ar[r]^{f} \ar[d]_{p}
  & G' \ar[d]^{p'} \\
G/H \ar[r]_{\overline{f}}
  & G'/H'\,.
}
\]
\end{lemma}

\begin{proof}
The first part is proved by reduction to the set-theoretic case. The second
follows immediately from~(i).
\end{proof}

One could carry over into the present situation the notions introduced above
for general equivalence relations.
% typo: équivalences -> équivalence
Let us merely point out the following lemma, whose proof is immediate by
reduction to the set-theoretic case:

\begin{lemma}{3.2.4}
\label{IV.3.2.4}
Let \(u\colon H\to G\) be a monomorphism of \(\mathcal{C}\)-groups and \(G'\)
a \(\mathcal{C}\)-subgroup of~\(G\). For the subobject \(G'\) of~\(G\) to be
stable under the equivalence relation defined by~\(H\), it is necessary and
sufficient that \(u\) factor through the canonical monomorphism \(G'\to G\),
and under this condition the operation of~\(H\) on~\(G'\) induced by the
operation of~\(H\) on~\(G\) defined by~\(u\) is none other than the operation
deduced from the monomorphism \(H\to G'\) through which \(u\) factors.
\end{lemma}

\par\addvspace{0.55\baselineskip}
\noindent\textbf{3.3. Universal effective equivalence relations}\label{IV.sec.3.3}
\par\addvspace{0.35\baselineskip}

\begin{definition}{3.3.1}
\label{IV.3.3.1}
Let \(f\colon X\to Y\) be a morphism. One calls the \emph{equivalence relation
defined by~\(f\) in~\(X\)}, and one denotes by \(\mathcal{R}(f)\), the
\(\widehat{\mathcal{C}}\)-equivalence relation in~\(X\) which is the image of
the canonical monomorphism
\[
X\times_{Y}X\longrightarrow X\times X.
\]
\end{definition}

\begin{definition}{3.3.2}
\label{IV.3.3.2}
Let \(R\) be an equivalence relation in~\(X\). One says that \(R\) is
\emph{effective} if
\begin{enumeratei}
\item \(R\) is representable (\ie is a \(\mathcal{C}\)-equivalence relation);
\item the quotient \(Y=X/R\) exists in~\(\mathcal{C}\)%
{\renewcommand{\thefootnote}{(13)}\nde{: One has added ``in \(\mathcal{C}\)''.}};
\item the diagram
\[
\xymatrix@C=3.6em{
R \ar@<0.55ex>[r]^{p_{1}} \ar@<-0.55ex>[r]_{p_{2}}
  & X \ar[r]^{p}
  & Y
}
\]
makes \(R\) the fiber square of~\(X\) over~\(Y\), that is, \(R\) is the
equivalence relation defined by~\(p\).
\end{enumeratei}
\end{definition}

\begin{scholium}{3.3.2.1}
\label{IV.3.3.2.1}
{\renewcommand{\thefootnote}{(14)}\footnote{\normalfont\textbf{N.D.E.}
One has added the numbering 3.3.2.1 (\resp~3.3.3.1) for later references.
On the other hand, one has added remark~3.3.3.2.}}%
\ If \(R\) is an effective equivalence relation in~\(X\), then \(p\) is an
effective epimorphism (1.3). If \(f\colon X\to Y\) is an effective
epimorphism, then \(\mathcal{R}(f)\) is an effective equivalence relation
in~\(X\), one quotient of which is~\(Y\). There is therefore a one-to-one
``Galois'' correspondence between effective equivalence relations in~\(X\)
and effective quotients of~\(X\) (\ie equivalence classes of effective
epimorphisms with source~\(X\)).
\end{scholium}

\begin{definition}{3.3.3}
\label{IV.3.3.3}
One says that the equivalence relation \(R\) in~\(X\) is \emph{universal
effective} if the quotient \(Y=X/R\) exists, and if, for every \(Y'\to Y\),
the fiber products \(X'=X\times_{Y}Y'\) and \(R'=R\times_{Y}Y'\) exist and
\(R'\) is a fiber square of~\(X'\) over~\(Y'\). It amounts to the same to say
that \(R\) is \emph{effective} and that \(p\colon X\to X/R\) is a
\emph{universal effective epimorphism}.
\end{definition}

\begin{scholium}{3.3.3.1}
\label{IV.3.3.3.1}
\textsuperscript{(14)}
There is therefore, as above, a one-to-one correspondence between universal
effective equivalence relations in~\(X\) and universal effective quotients
of~\(X\).
\end{scholium}

\begin{remark}{3.3.3.2}
\label{IV.3.3.3.2}
\textsuperscript{(14)}
Suppose that \(\mathcal{C}\) is the category of \(S\)-schemes, and denote
by~\(\mathbf{A}^{1}\) the affine space of dimension~\(1\) over~\(S\). Let
\(R\subset X\times_{S}X\) be a universal effective equivalence relation and
\(p\colon X\to Y\) the quotient. Then, for every open~\(U\) of~\(Y\),
\(\mathcal{O}(U)=\Hom_{S}(U,\mathbf{A}^{1}_{S})\) is the set of elements
\(\phi\) of
\(\mathcal{O}(p^{-1}(U))=\Hom_{S}(p^{-1}(U),\mathbf{A}^{1}_{S})\) such that
\(\phi\circ\mathrm{pr}_{1}=\phi\circ\mathrm{pr}_{2}\). In particular, if
\(R\) is given by the action of a group \(H\) operating freely on the right
on~\(X\) (\cf~3.2.1), then \(\mathcal{O}(U)\) is the set of
\(\phi\in\mathcal{O}(p^{-1}(U))\) such that \(\phi(xh)=\phi(x)\), for every
\(S'\to S\) and \(x\in X(S')\), \(h\in H(S')\).
\end{remark}

\begin{proposition}{3.3.4}
\label{IV.3.3.4}
Let \(R\) be a universal effective equivalence relation in~\(X\). Let
\(f\colon X\to Z\) be a morphism compatible with~\(R\), hence factoring
through \(g\colon X/R\to Z\). The following conditions are equivalent:
\begin{enumeratei}
\item \(g\) is a monomorphism;
\item \(R\) is the equivalence relation defined by~\(f\).
\end{enumeratei}
\end{proposition}

\begin{proof}
Indeed, (i)~implies~(ii) trivially, and the converse is none other than~2.5.
\end{proof}

\begin{definition}{3.3.5}
\label{IV.3.3.5}
Let \(H\) be a \(\mathcal{C}\)-group operating freely on~\(X\). One says that
\(H\) operates in an \emph{effective} manner, or that the operation of~\(H\)
on~\(X\) is \emph{effective} (\resp\ \emph{universal effective}), if the
equivalence relation defined in~\(X\) by the action of~\(H\) is effective
(\resp\ universal effective).
\end{definition}

\par\addvspace{0.55\baselineskip}
\noindent\textbf{3.4. \((M)\)-effectivity.}\ ---
\label{IV.sec.3.4}
In practice, it is most often difficult to characterize the universal
effective epimorphisms. One nevertheless often has at one's disposal a
certain number of morphisms of this type, for example in the theory of
schemes, the faithfully flat quasi-compact morphisms. This leads to the
developments below.

\numpara{3.4.1}
\label{IV.3.4.1}
Let us first state a certain number of conditions on a family~\((M)\) of
morphisms of~\(\mathcal{C}\):
\begin{enumeratea}
\item \emph{\((M)\) is stable under base extension, \ie every
\(u\colon T\to S\) which is an element of~\((M)\) is quarrable
(\cf~1.4.0) and for every \(S'\to S\),
\(u'\colon T\times_{S}S'\to S'\) is an element of~\((M)\).}
\item \emph{The composite of two elements of~\((M)\) lies in~\((M)\).}
\item \emph{Every isomorphism is an element of~\((M)\).}
\item \emph{Every element of~\((M)\) is an effective epimorphism.}
\end{enumeratea}

Let us note that~(a) and~(b) imply:

\emph{(a') The cartesian product of two elements of~\((M)\) lies
in~\((M)\): Let \(u\colon X\to Y\) and \(u'\colon X'\to Y'\) be two
\(S\)-morphisms which are elements of~\((M)\). If \(Y\times_{S}Y'\) exists,
then \(X\times_{S}X'\) exists and \(u\times_{S}u'\) is an element
of~\((M)\).}

This follows from the diagram
\[
\xymatrix@C=2.6em@R=2.4em{
  & Y'
  & X' \ar[l]_{u'}
\\
X \ar[d]_{u}
  & X\times_{S}Y' \ar[l] \ar[u] \ar[d]
  & X\times_{S}X' \ar[l] \ar[u] \ar[dl]^{u\times_{S}u'}
\\
Y
  & Y\times_{S}Y' \ar[l]
  &
}
\,.
\]
Likewise~(a) and~(d) imply:

\emph{(d') Every element of~\((M)\) is a universal effective epimorphism.}

\numpara{3.4.2}
\label{IV.3.4.2}
The family \((M_{0})\) of universal effective epimorphisms satisfies
conditions~(a) to~(d) of~3.4.1. Indeed, (a), (c) and~(d) are satisfied by
definition, and~(b) follows from~1.8. In what follows, we shall assume given
a family~\((M)\) of morphisms of~\(\mathcal{C}\) satisfying these conditions:
our results will therefore apply to the family~\((M_{0})\) in particular.

\begin{definition}{3.4.3}
\label{IV.3.4.3}
One says that the equivalence relation \(R\) in~\(X\) is \emph{of type~\((M)\)}
if it is representable and if \(p_{1}\in(M)\) (which by~(b) and~(c) implies
that \(p_{2}\in(M)\)).

One says that \(R\) is \emph{\((M)\)-effective} if it is effective and if the
canonical morphism \(X\to X/R\) is an element of~\((M)\).

One says that the quotient \(Y\) of~\(X\) is \emph{\((M)\)-effective} if the
canonical morphism \(X\to Y\) is an element of~\((M)\).
\end{definition}

From this definition there follow the following consequences:%
{\renewcommand{\thefootnote}{(15)}\nde{: One has added the numbering
3.4.3.1, for later references, and one has set out in detail the proof of
point~(i).}}

\begin{proposition}{3.4.3.1}
\label{IV.3.4.3.1}
\begin{enumeratei}
\item An \((M)\)-effective equivalence relation is of type~\((M)\) and
universal effective.
\item An \((M)\)-effective quotient is universal effective (\cf~3.3.3).
\item The maps \(R\mapsto X/R\) and \(p\mapsto\mathcal{R}(p)\) realize a
one-to-one correspondence between \((M)\)-effective equivalence relations
in~\(X\) and \((M)\)-effective quotients of~\(X\).
\item \((M_{0})\)-effective is equivalent to universal effective.
\end{enumeratei}
\end{proposition}

\begin{proof}
Let us prove point~(i). Since \(R\) is \((M)\)-effective, one has a cartesian
square
\[
\xymatrix@C=4.2em@R=2.4em{
R \ar[r]^{p_{2}} \ar[d]_{p_{1}}
  & X \ar[d]^{p} \\
X \ar[r]_{p}
  & X/R\,,
}
\]
and \(p\in(M)\). Then, by~3.4.1~(a), \(p_{1}\) and \(p_{2}\) belong
to~\((M)\), hence \(R\) is of type~\((M)\).

Set \(Y=X/R\), and let \(Y'\to Y\) be an arbitrary morphism. By~3.4.1~(a),
the fiber products \(X'=X\times_{Y}Y'\) and \(R'=R\times_{Y}Y'\) exist, and
the morphisms \(X'\to Y'\) and \(p'_{i}\colon R'\to X'\) belong to~\((M)\).
Finally, since \(R=X\times_{Y}X\), one obtains, by associativity of the
fiber product:
\[
R'=X\times_{Y}X\times_{Y}Y'=X'\times_{Y'}X'.
\]
This shows that \(R'\) is \((M)\)-effective; hence, in particular, \(R\) is
universal effective. This proves~(i), and also~(iv). Points~(ii) and~(iii)
follow from this, taking definition~3.3.2 into account.
\end{proof}

\numpara{3.4.4}
\label{IV.3.4.4}
Let \(H\) be an \(S\)-group whose structural morphism is an element
of~\((M)\). Then if \(H\) operates \emph{freely} on the \(S\)-object~\(X\), it
defines on it an equivalence relation of type~\((M)\). Indeed, by~(a) the
fiber product \(H\times_{S}X\) exists and
\(\mathrm{pr}_{2}\colon H\times_{S}X\to X\) is an element of~\((M)\). One
will say that the operation of~\(H\) is \emph{\((M)\)-effective} if the
equivalence relation defined in~\(X\) by this operation is
\((M)\)-effective.

\begin{proposition}{3.4.5 ($(M)$-effectivity and base change)}
\label{IV.3.4.5}
Let \(R\) be an \((M)\)-effective equivalence relation in~\(X\) over~\(S\).
Set \(Y=X/R\). Let \(S'\to S\) be a base change such that
\(Y'=Y\times_{S}S'\) exists. Then \(X'=X\times_{S}S'\) exists,
\(R'=R\times_{S}S'\) is an \((M)\)-effective equivalence relation in~\(X'\)
over~\(S'\), and \(X'/R'\simeq(X/R)'\).
\end{proposition}

\begin{proof}
Indeed, the canonical morphisms \(X\to Y\) and \(R\to Y\) are elements
of~\((M)\), hence by~(a') \(X'\) and \(R'\) are representable. By
associativity of the product, \(R'\) is the equivalence relation defined
in~\(X'\) by the canonical morphism \(X'\to Y'\), which is an element
of~\((M)\), whence the conclusion.
\end{proof}

\begin{proposition}{3.4.6 ($(M)$-effectivity and cartesian products)}
\label{IV.3.4.6}
Let \(R\) (\resp~\(R'\)) be an \((M)\)-effective equivalence relation in~\(X\)
(\resp~\(X'\)) over~\(S\). If \((X/R)\times_{S}(X'/R')\) exists, then
\(X\times_{S}X'\) exists, \(R\times_{S}R'\) is an \((M)\)-effective
equivalence relation in \(X\times_{S}X'\) over~\(S\), and
\[
(X\times_{S}X')/(R\times_{S}R')
\simeq
(X/R)\times_{S}(X'/R').
\]
\end{proposition}

\begin{proof}
Set \(Y=X/R\), \(Y'=X'/R'\). By~(a'), the fiber product
\(X\times_{S}X'\) exists and the canonical morphism
\[
q\colon X\times_{S}X'\longrightarrow Y\times_{S}Y'
\]
is an element of~\((M)\). Now the formula
\[
(X\times X')\times_{(Y\times Y')}(X\times X')
\simeq
(X\times_{Y}X)\times(X'\times_{Y'}X')
\]
(all unindexed products are taken over~\(S\)) shows that \(R\times_{S}R'\)
is the equivalence relation defined by~\(q\) in \(X\times_{S}X'\), which
completes the proof.
\end{proof}
